\subsection{Adapter: Safe Copies and Demand-Driven Recovery}
\label{sec:kv-runtime}

\paragraph{Independent copies and publication.}
Adapter (4) applies the Planner's decisions through the Connector, which
performs allocation, transfer, and completion notification. It does not
choose tiers, destinations, or hot sets. Peer and fault-independent DRAM
copies have separate allocation, transfer, and commit paths. Failure of
one path does not invalidate the other's committed objects. A Peer copy
that is the only valid surviving replica cannot be discarded before
DRAM commit. Shared prefixes occupy one physical copy and can share
the ordinary prefix cache's DRAM representation.

A proposed placement is not a readable replica. Publication follows
transfer completion and validation of model, layout, and shard identity.
Selector state is copied with the initial layer and made available before
selection. Bytes in copying, committed, or reclaiming states count
toward physical occupancy until release is complete.

\paragraph{Foreground protection.}
When foreground KV exceeds the watermark, the Adapter reclaims Peer
footprints in reverse admission priority: larger-slack requests first,
and later layers first within a footprint being removed. The policy
withdraws a request's complete footprint; the layer order controls the
physical release sequence. The request competes again at the next window.
DRAM coverage makes this downgrade safe. Low-priority copy streams are
rate-limited (default: 20\% of link capacity), with target HBM bandwidth
also accounted for. For scale, a 64-token seal across all MiniCPM4 layers
produces 2\,MiB; 100 requests at 50 tokens/s generate about 160\,MiB/s
of new KV. These are byte-rate calculations, not throughput measurements.

\paragraph{Per-request loading and readiness.}
For each retry the Adapter preloads by increasing layer and decreasing
hotness, skipping locally resident Peer blocks. When a selector reveals
exact demand, missing blocks receive higher priority than remaining
preloads. The analytic model assumes preemption; whether an in-flight
DMA can actually be preempted depends on the transfer implementation.
Once a request's needed KV for layer $l$ is ready, that layer may compute
without waiting for another request. Later preloads overlap current
computation. Cold KV is filled in the background, and subsequent decode
can demand-fetch still-absent blocks. Tier R recomputes the history.
No retry waits solely for all historical KV to arrive.

\paragraph{Recovery invariant.}
AegisKV changes placement and loading order, not the selector's output.
Every required attention input must be available before the corresponding
attention operation executes. With the same weights, compatible layout
and operators, and the known token prefix, recovery reconstructs the
same logical KV state without adding an attention approximation.
Floating-point ordering can still change with batching; this mechanism
does not by itself guarantee token-identical future stochastic sampling.
